\subsection{同态模中基环的扩张}

设 A 和 B 是两个环，E 是右 A-模，F 是右 B-模，G 是 (B, A)-双模。回顾我们曾定义（《代数学》第二章 §4 第 2 目）一个典范 Z-模同态

\[
\nu \colon \mathbf {F} \otimes_ {\mathbf {B}} \operatorname{Hom} _ {\mathbf {A}} (\mathbf {E}, \mathbf {G}) \rightarrow \operatorname{Hom} _ {\mathbf {A}} (\mathbf {E}, \mathbf {F} \otimes_ {\mathbf {B}} \mathbf {G})\tag{8}
\]

它使得对所有 \(y \in \mathbf{F}\) 和 \(u \in \operatorname{Hom}_{\mathbf{A}}(\mathbf{E}, \mathbf{G})\) ， \(\nu(y \otimes u)\) 是 A-线性映射 \(x \mapsto y \otimes u(x)\) 。

命题 10. 设 A、B 是两个环，E 是右 A-模，F 是右 B-模，G 是 (B, A)-双模。设 F 是平坦的。则若 E 是有限生成的（相应地有限表示的），典范同态 (8) 是单射的（相应地双射的）。

把 A、B、F、G 视为固定，并对每个右 A-模 E 令

\[
\mathrm{T} (\mathrm{E}) = \mathrm{F} \otimes_ {\mathrm{B}} \operatorname{Hom} _ {\mathrm{A}} (\mathrm{E}, \mathrm{G}), \quad \mathrm{T} ^ {\prime} (\mathrm{E}) = \operatorname{Hom} _ {\mathrm{A}} (\mathrm{E}, \mathrm{F} \otimes_ {\mathrm{B}} \mathrm{G})
\]

且用 \(\nu_{\mathbf{E}}\) 表示同态 (8)；对每个右 A-模同态 \(\nu: \mathbf{E} \to \mathbf{E}'\) ，令 \(\mathrm{T}(v) = 1_{\mathbf{F}} \otimes \mathrm{Hom}(v, 1_{\mathbf{G}})\) 和 \(\mathrm{T}'(v) = \mathrm{Hom}(v, 1_{\mathbf{F}} \otimes 1_{\mathbf{G}})\) 。设 \(\mathrm{L}_1 \xrightarrow{v} \mathrm{L}_0 \xrightarrow{w} \mathrm{E} \to 0\) 是 E 的一个表示；我们假设自由模 \(\mathrm{L}_0\) （相应地自由模 \(\mathrm{L}_0\) 和 \(\mathrm{L}_1\) ）是有限生成的。图

\[
\begin{array}{c} 0 \longrightarrow \mathrm{T(E)} \xrightarrow {\mathrm{T(w)}} \mathrm {T(L_ {0})} \xrightarrow {\mathrm{T(v)}} \mathrm {T(L_ {1})} \\ v _ {\mathbf {E}} \Biggl \downarrow \qquad \qquad \qquad \qquad \qquad \qquad \nu_ {\mathrm {L_ {0}}} \Biggl \downarrow \qquad \qquad \qquad \nu_ {\mathrm {L_ {1}}} \Biggl \downarrow \\ 0 \longrightarrow \mathrm {T^ {\prime} (E)} \xrightarrow {\mathrm {T^ {\prime} (w)}} \mathrm {T^ {\prime} (L_ {0})} \xrightarrow {\mathrm {T^ {\prime} (v)}} \mathrm {T^ {\prime} (L_ {1})} \end{array}\tag{9}
\]

是交换的，且第二行是正合的（《代数学》第二章 §2 第 1 目定理 1）；此外，序列

\[
0 \rightarrow \operatorname{Hom} _ {\mathbf {A}} (\mathrm{E}, \mathrm{G}) \rightarrow \operatorname{Hom} _ {\mathbf {A}} (\mathrm{L} _ {0}, \mathrm{G}) \rightarrow \operatorname{Hom} _ {\mathbf {A}} (\mathrm{L} _ {1}, \mathrm{G})
\]

是正合的（出处同上），且由于 F 是平坦的，(9) 的第一行也是一个正合列（第 3 目，命题 1）。于是我们知道 \(\nu_{\mathrm{L}_0}\) （相应地 \(\nu_{\mathrm{L}_0}\) 和 \(\nu_{\mathrm{L}_1}\) ）是双射的（《代数学》第二章 §4 第 2 目命题 2）。若只假设 \(\nu_{\mathrm{L}_0}\) 是双射的，则由 (9) 得

\[
\nu_ {\mathbf {L} _ {0}} \circ \mathbf {T} (w) = \mathbf {T} ^ {\prime} (w) \circ \nu_ {\mathbf {E}}
\]

是单射的，因而 \(v_{E}\) 也是单射的。若假设 \(v_{L_{0}}\) 和 \(v_{L_{1}}\) 都是双射的，则由 §1 第 4 目命题 2 的推论 2 (ii) 得 \(v_{E}\) 是满射的，而由于我们刚看到 \(v_{E}\) 是单射的，它是双射的。

